\documentclass[11pt]{amsart}

\usepackage[a4paper,margin=1.02in]{geometry}
\usepackage{amsmath,amssymb,amsthm,mathtools,mathrsfs,bm}
\usepackage{booktabs,tabularx,array,longtable}
\usepackage{enumitem}
\usepackage{tikz-cd}
\usepackage{microtype}
\usepackage{xcolor}
\usepackage{hyperref}
\usepackage[nameinlink,capitalize,noabbrev]{cleveref}
\usepackage[most]{tcolorbox}
\usepackage{listings}

\hypersetup{
  colorlinks=true,
  linkcolor=blue!45!black,
  citecolor=blue!45!black,
  urlcolor=blue!45!black,
  pdftitle={Decisive Obstacles in the IUT--ABC Formalization Programme},
  pdfauthor={ChatGPT}
}

\newtheorem{theorem}{Theorem}[section]
\newtheorem{proposition}[theorem]{Proposition}
\newtheorem{lemma}[theorem]{Lemma}
\newtheorem{corollary}[theorem]{Corollary}
\newtheorem{counterexample}[theorem]{Counterexample}
\newtheorem{criterion}[theorem]{Criterion}
\theoremstyle{definition}
\newtheorem{definition}[theorem]{Definition}
\newtheorem{construction}[theorem]{Construction}
\theoremstyle{remark}
\newtheorem{remark}[theorem]{Remark}
\newtheorem{status}[theorem]{Formalization status}

\newcommand{\R}{\mathbb R}
\newcommand{\Q}{\mathbb Q}
\newcommand{\N}{\mathbb N}
\newcommand{\Z}{\mathbb Z}
\newcommand{\OO}{\mathcal O}
\newcommand{\eps}{\varepsilon}
\newcommand{\rad}{\operatorname{rad}}
\newcommand{\Gal}{\operatorname{Gal}}
\newcommand{\Spec}{\operatorname{Spec}}
\newcommand{\im}{\operatorname{im}}
\newcommand{\Vol}{\operatorname{Vol}}
\newcommand{\Corr}{\operatorname{Corr}}
\newcommand{\Diff}{\operatorname{Diff}}
\newcommand{\Err}{\operatorname{Err}}
\newcommand{\Cond}{\operatorname{Cond}}
\newcommand{\ThetaOut}{\mathcal U_{\Theta}}
\newcommand{\ThetaEnv}{\mathcal E_{\Theta}}
\newcommand{\thetaK}{\vartheta}

\title[Decisive IUT--ABC obstacles]{Decisive Obstacles in the IUT--$abc$ Formalization Programme\\
\large Local theta-root geometry, relational multiradial source realization, and quantitative auxiliary-prime control}
\author{ChatGPT}
\date{August 24, 2026}

\subjclass[2020]{Primary 11D75, 11G50; Secondary 11R58, 14G22, 03B35, 28C10}
\keywords{ABC conjecture, inter-universal Teichm\"uller theory, Lean, Tate curve, Kummer cover, Chebyshev function, log-volume, formal verification}

\begin{document}
\maketitle

\begin{abstract}
We isolate and attack the three remaining bottlenecks closest to the final
IUT--$abc$ implication in the accompanying Lean development.  The first is
the construction of an actual anabelian/tempered initial-theta fiber.  We
separate its automatic choices from its genuine geometric content and reduce
the local theta-root model to separable Kummer fibers
$Y^{\ell}-\thetaK_q(u)$ away from the Tate theta zero orbit.  The second is the
realization of the source of IUT III, Theorem~3.11.  We give a
source-faithful relational formalism in which Ind1 and Ind2 are deterministic
transformations while Ind3 is an upper-semicompatible relation, and prove that
a pointwise norm estimate for the log-link is sufficient to control the full
possible-image union.  The third is the quantifier problem in the auxiliary
prime.  We prove two independent alternatives: a bounded-interval
Chebyshev-mass escape criterion, and an asymptotic relative-defect criterion
under which no upper bound for the selected prime is required.  We also record
counterexamples showing that none of the three relative defects may be ignored
in isolation.

The paper distinguishes ordinary mathematical proofs, Lean-kernel-verified
files, files undergoing continuous integration, and genuine source theorems
that remain unproved.  The resulting dependency graph is non-circular and
ends in the standard logarithmic $abc$ conjecture once the stated geometric
and uniform asymptotic source theorems are supplied.  No unconditional proof
of $abc$ is claimed.
\end{abstract}

\begin{tcolorbox}[colback=yellow!5,colframe=black!65,title=Honesty boundary]
The verified downstream chain proves implications of the form
\[
 \text{actual initial-theta family}
 \Longrightarrow
 \text{actual IUT III source}
 \Longrightarrow
 \text{IUT IV estimate}
 \Longrightarrow
 abc.
\]
This paper improves the first three antecedents and proves several exact
reduction theorems.  It does not construct every antecedent unconditionally.
In particular, there is presently no parameter-free Lean declaration
\[
 \texttt{theorem abc\_conjecture : ABCConjecture}.
\]
\end{tcolorbox}

\tableofcontents

\section{Target statement and proof architecture}

For a positive integer $n$, write
\[
 \rad(n)=\prod_{p\mid n}p.
\]
The target is the logarithmic $abc$ conjecture.

\begin{theorem}[Target]\label{thm:abc-target}
For every $\eps>0$ there exists $C_{\eps}\in\R$ such that every positive,
pairwise coprime triple $a+b=c$ satisfies
\[
 \log \max\{a,b,c\}
 \le (1+\eps)\log\rad(abc)+C_{\eps}.
\]
\end{theorem}

The formal programme is organized into the following dependency chain:
\[
\begin{tikzcd}[column sep=small]
\text{Frey arithmetic}
 \arrow[r]
&\text{initial-theta data}
 \arrow[r]
&\text{Theorem 3.11 source}
 \arrow[r]
&\text{Corollary 3.12}
 \arrow[r]
&\text{IUT IV q-bound}
 \arrow[r]
&abc.
\end{tikzcd}
\]
The elementary arithmetic, the scalar IUT IV algebra, and the final
height--radical passage have been substantially formalized.  The decisive
open constructions lie in the middle of the chain.

\section{Research-group decomposition}

To avoid premature convergence, the investigation is divided into independent
cells.

\begin{enumerate}[label=\textbf{Cell \Alph*:},leftmargin=2.7em]
\item \emph{Fiber geometry}: construct actual local theta-root covers,
orbicurves, cores, cartesian diagrams, graph cusps, and tempered comparison
maps.
\item \emph{Multiradial source}: realize ordinary branches and Ind1--Ind3
operations from genuine Hodge theaters and prove their pointwise norm or
volume control.
\item \emph{Auxiliary prime and uniformity}: pursue both bounded-prime
selection and unbounded-prime relative-error absorption.
\item \emph{Adversarial audit}: construct countermodels to overstrong
interfaces and identify quantifier orders that cannot be recovered from
weaker hypotheses.
\item \emph{Formal verification}: maintain a separate Lean branch for each
mathematical direction and merge only after a pinned kernel build succeeds.
\end{enumerate}

A direction is not discarded merely because it lacks a current library
implementation.  It is removed only after a strict counterexample,
contradiction, or theorem showing that the proposed implication is impossible.

\section{The Frey local input}

For an $abc$ point $P=(a,b,c)$, consider the integral Frey model
\[
 E_P:\qquad y^2=x(x-a)(x+b).
\]
Its standard invariants are
\[
 c_4(E_P)=16(a^2+ab+b^2),
 \qquad
 \Delta(E_P)=16(abc)^2.
\]
The quadratic core is coprime to $abc$ for a primitive triple.

\begin{proposition}[Odd multiplicative signature]\label{prop:odd-signature}
Let $p$ be an odd prime dividing $abc$.  Then
\[
 p\mid \Delta(E_P),
 \qquad
 p\nmid c_4(E_P).
\]
Moreover, every primitive positive $abc$ point other than $(1,1,2)$ admits
such an odd prime.
\end{proposition}

\begin{proof}
The divisibility $p\mid\Delta(E_P)$ follows immediately from
$\Delta(E_P)=16(abc)^2$.  Since $p$ is odd, it does not divide $16$.  If
$p\mid c_4(E_P)$, then $p$ divides $a^2+ab+b^2$.  But the primitive relations
and $a+b=c$ imply
\[
 \gcd\bigl(abc,a^2+ab+b^2\bigr)=1,
\]
a contradiction.

If no odd prime divides $abc$, then each of $a,b,c$ is a power of $2$.
Pairwise coprimality implies at most one of them is even.  Since $a+b=c$ and
$a,b>0$, the only possibility is $a=b=1$ and $c=2$.
\end{proof}

\begin{lemma}[Unit $c_4$ forces minimality]\label{lem:unit-c4-minimal}
Let $R$ be a discrete valuation ring with fraction field $K$.  Let $W$ be an
integral Weierstrass equation over $R$.  If
\[
 v(c_4(W))=1
\]
in the multiplicative valuation convention, then $W$ is minimal.
\end{lemma}

\begin{proof}
For an admissible variable change $C$ with scaling parameter $u$, one has
\[
 c_4(C\cdot W)=u^{-4}c_4(W),
 \qquad
 \Delta(C\cdot W)=u^{-12}\Delta(W).
\]
Assume $C\cdot W$ remains integral.  Integrality gives
$v(c_4(C\cdot W))\le 1$.  Hence
\[
 v(u^{-1})^4v(c_4(W))\le1.
\]
Since $v(c_4(W))=1$, we obtain $v(u^{-1})^4\le1$, and therefore
$v(u^{-1})\le1$.  It follows that
\[
 v(\Delta(C\cdot W))
 =v(u^{-1})^{12}v(\Delta(W))
 \le v(\Delta(W)).
\]
Thus the original integral equation has maximal discriminant valuation among
all integral equations in its isomorphism class, which is the definition of
minimality used in the formal library.
\end{proof}

\begin{corollary}[Local multiplicative reduction criterion]
Under the hypotheses of \cref{lem:unit-c4-minimal}, if in addition
\[
 v(\Delta(W))<1,
\]
then $W$ has multiplicative reduction.
\end{corollary}

\begin{status}
The odd-support and exact invariant calculations are Lean-kernel verified on
\texttt{main}.  The generic constructor for multiplicative reduction is also
verified.  The direct proof of \cref{lem:unit-c4-minimal} is maintained on a
separate branch until its pinned build succeeds.  The arithmetic transfer from
integer divisibility to maximal-ideal membership in the completed local ring
remains a separate theorem.
\end{status}

\section{Actual global core: what has and has not been constructed}

The formal development now proves a precise assembler theorem.  Given a
number field $F$, an algebraic closure $\overline F$, full rationality of
$E_P[12]$, a square root of $-1$, stable reduction away from $2$ and away from
$3$, and one odd multiplicative bad moduli place, one obtains the public global
initial-theta conditions and a fixed IUT core.

\begin{theorem}[Fixed-core assembly]\label{thm:fixed-core-assembly}
An explicit twelve-torsion global input package for the Frey curve determines
an actual fixed global core, calibrated by
\[
 j(E_P/F)=\operatorname{algebraMap}_{\Q,F}(j(E_P/\Q)).
\]
\end{theorem}

\begin{proof}[Proof sketch]
The inclusion $E[6]\subseteq E[12]$ follows from $6\mid12$, so full
$12$-torsion rationality implies the public six-torsion condition.  Stable
reduction away from $2$ and away from $3$ covers every finite place: a residue
characteristic cannot equal both $2$ and $3$.  Taking the singleton containing
the supplied odd multiplicative moduli place gives a nonempty bad locus.  The
remaining structure fields are then assembled definitionally.
\end{proof}

This is an assembly theorem, not an unconditional existence theorem.  The
remaining standard sources are the Weil-pairing/cyclotomic input, the
semistable-reduction theorem at levels $3$ and $4$, and the completed local
multiplicative place supplied by \cref{prop:odd-signature} and the local
valuation transfer.

\section{Removing automatic choices from the fiber gap}

Several fields in the public interface are not independent geometric
obstructions.

\begin{proposition}[Canonical quotient element and cusp]
Suppose a rank-one quotient is identified with $\Z/\ell\Z$ by
\[
 \phi:Q\overset\sim\longrightarrow \Z/\ell\Z,
\]
where $\ell$ is prime.  Then
\[
 q:=\phi^{-1}(1)
\]
is nonzero, and the distinguished cusp may be defined as the cusp associated
to $q$.
\end{proposition}

\begin{proof}
Since $\ell$ is prime, $1\ne0$ in $\Z/\ell\Z$.  Hence
$\phi(q)=1\ne0$, so $q\ne0$.  The cusp specification is then definitional.
\end{proof}

\begin{proposition}[Presently weak decomposition-group choice]
If the public local interface asks only for a closed subgroup and states no
place-stabilizer characterization, the full Galois group is a canonical closed
choice.
\end{proposition}

\begin{remark}
This proposition does not identify the full group with the genuine
decomposition group.  It proves that the current interface is too weak to
express that arithmetic meaning.  A faithful interface should add a
characterizing condition such as
\[
 \sigma\in D_w\quad\Longleftrightarrow\quad \sigma(w)=w,
\]
up to the intended conjugacy convention.
\end{remark}

Consequently, the genuine fiber gap consists of the orbicurves and covers,
type and core theorems, cartesian diagrams, corrected valuation section,
theta-root model, tempered comparison, and graph-cusp compatibility.  The
quotient representative, its nonzero proof, the derived cusp, and the weak
closed-subgroup field are automatic.

\section{Pointwise theta Kummer fibers}

Let $K$ be a complete nonarchimedean normed field, let $q\in K$ satisfy
$0<|q|<1$, and let $\thetaK_q(u)$ denote the Tate theta product.  The available
theta divisor theorem says that
\[
 \thetaK_q(u)=0
 \quad\Longleftrightarrow\quad
 u=-q^k\text{ for some }k\in\Z.
\]

\begin{definition}[Theta Kummer polynomial]
For $\ell\ge1$ and $u\in K^\times$, define
\[
 f_{q,u,\ell}(Y)=Y^\ell-\thetaK_q(u)\in K[Y].
\]
\end{definition}

\begin{theorem}[Separable theta Kummer fiber]\label{thm:kummer-fiber}
Assume $u\ne-q^k$ for all $k\in\Z$, and assume $\ell$ is invertible in $K$.
Then $f_{q,u,\ell}$ is monic and separable.  Every geometric root is simple.
The algebra
\[
 K[Y]/(f_{q,u,\ell})
\]
contains a canonical element $y$ satisfying
\[
 y^\ell=\thetaK_q(u),
 \qquad y\ne0.
\]
\end{theorem}

\begin{proof}
The zero-orbit hypothesis gives $\thetaK_q(u)\ne0$.  The polynomial is monic
because its leading term is $Y^\ell$.  Its derivative is
\[
 f'_{q,u,\ell}(Y)=\ell Y^{\ell-1}.
\]
If a geometric element $x$ were a common root of $f$ and $f'$, invertibility
of $\ell$ would imply $x=0$.  But then
$f(0)=-\thetaK_q(u)\ne0$, a contradiction.  Thus $f$ is separable and all its
roots are simple.  The canonical class of $Y$ in the quotient satisfies the
power equation by construction.  If it were zero, then its positive
$\ell$-th power would be zero, contradicting
$\thetaK_q(u)\ne0$.
\end{proof}

\begin{status}
The pointwise Kummer construction is encoded on a research branch using the
actual theta-product divisor theorem and the library's Kummer-polynomial API.
The remaining geometric task is to globalize these fibers over
$K^\times/q^{\Z}$, prove period descent and ramification behavior, and identify
the resulting cover with the required orbicurve and tempered object.
\end{status}

\section{A relational model of the genuine Theorem 3.11 source}

An exact word normal form for every possible image is stronger than the source
statement needed in IUT III.  The Ind3 ambiguity is naturally relational.

\begin{definition}[Upper-semicompatible source system]
Fix a carrier $X$.  A source system consists of:
\begin{itemize}
\item ordinary regions $U_o\subseteq X$;
\item deterministic operations $T_{1,a}$ and $T_{2,b}$ on regions;
\item a relation $R_{3,c}(U,V)$ for Ind3;
\item a distinguished ordinary region $U_{\mathrm{nat}}$;
\item an envelope $\ThetaEnv\subseteq X$.
\end{itemize}
Reachability is the least relation generated by ordinary regions, Ind1,
Ind2, and all relational steps $R_{3,c}(U,V)$.
\end{definition}

\begin{theorem}[Relational envelope theorem]\label{thm:relational-envelope}
Assume:
\begin{enumerate}[label=(\roman*)]
\item every ordinary region lies in $\ThetaEnv$;
\item Ind1 preserves containment in $\ThetaEnv$;
\item Ind2 preserves containment in $\ThetaEnv$;
\item whenever $U\subseteq\ThetaEnv$ and $R_{3,c}(U,V)$, one has
$V\subseteq\ThetaEnv$.
\end{enumerate}
Then
\[
 U_{\mathrm{nat}}\subseteq\ThetaOut
 \subseteq\ThetaEnv,
\]
where $\ThetaOut$ is the union of all reachable regions.
\end{theorem}

\begin{proof}
The native inclusion holds because the distinguished ordinary region is one
of the constructors of reachability.  For the envelope inclusion, perform
induction on the reachability derivation.  The ordinary case is (i), the Ind1
and Ind2 cases are (ii) and (iii), and the relational Ind3 case is (iv).
Taking the union over all reachable regions gives the result.
\end{proof}

The point of the relational formulation is that Ind3 may be verified by a
local map rather than an arbitrary set-level volume assertion.

\begin{theorem}[Pointwise norm reduction]\label{thm:norm-reduction}
Let $f:X\to Y$ be a map between seminormed additive groups satisfying
\[
 \|f(x)\|\le\|x\|\qquad\text{for all }x.
\]
Let
\[
 B_X(r)=\{x:\|x\|\le r\},
 \qquad
 B_Y(r)=\{y:\|y\|\le r\}.
\]
If $U\subseteq B_X(r)$ and $V\subseteq f(U)$, then
\[
 V\subseteq B_Y(r).
\]
The same assertion holds coordinatewise for a product packet.
\end{theorem}

\begin{proof}
For $y\in V$, choose $x\in U$ with $f(x)=y$.  Then
\[
 \|y\|=\|f(x)\|\le\|x\|\le r.
\]
The product statement follows by applying the same estimate in each
coordinate.
\end{proof}

\begin{corollary}[Source obligations reduced to local estimates]
If the genuine Ind1, Ind2, and log-link maps are pointwise
norm-nonincreasing, and the genuine Ind3 target is contained in the image of
the source under the log-link, then the complete possible-image union is
contained in the radial packet envelope.
\end{corollary}

Thus the set-theoretic portion of the genuine Theorem~3.11 source reduces to:
\begin{enumerate}[label=(\alph*)]
\item identify the actual ordinary branches;
\item prove norm control for procession reindexing;
\item prove norm control for spectral-factor automorphisms;
\item prove norm control for the log-link;
\item identify the public packet coordinates and the genuine measured hull.
\end{enumerate}

\section{The native volume and Corollary 3.12}

Once the source relations
\[
 U_{\mathrm{nat}}\subseteq\ThetaOut\subseteq\ThetaEnv
\]
are established in a finite-positive measured realization, monotonicity gives
\[
 \log\mu(U_{\mathrm{nat}})
 \le \log\mu(\ThetaOut)
 \le \log\mu(\ThetaEnv).
\]
The intended calibration is
\[
 \log\mu(U_{\mathrm{nat}})=-|\log q|.
\]
If the theta coefficient is defined by
\[
 -|\log\Theta|=C_{\Theta}|\log q|,
\]
then the native inclusion yields $C_{\Theta}\ge-1$.

The formal downstream implication is already available.  The genuine source
work is the construction of the Haar-measured packet, the actual hull or outer
measure, and the identification of the native packet volume with the
procession-normalized $q$-pilot value.

\section{Why eventual large image is not enough}

A Serre-style theorem normally states that a desired representation-theoretic
property holds for every sufficiently large prime.  Such an eventual statement
supplies no independent upper bound for the first usable prime.

\begin{counterexample}[Eventuality without a bounded witness]
Fix $B\in\N$ and define a prime property by
\[
 P(p)\quad\Longleftrightarrow\quad B<p.
\]
Then $P$ holds for every sufficiently large prime, but no prime satisfying
$P$ is at most $B$.
\end{counterexample}

Therefore a proof requiring both
\[
 p>N
 \qquad\text{and}\qquad
 p\le B
\]
cannot obtain the second inequality from eventuality alone.  One must either
prove a quantitative prime theorem or eliminate the need for an upper bound.

\section{Bounded-prime route via Chebyshev mass}

Let
\[
 \theta(X)=\sum_{p\le X}\log p
\]
be the Chebyshev theta function.  For a finite forbidden set $A$ of primes,
define
\[
 m(A)=\sum_{p\in A}\log p.
\]

\begin{theorem}[Chebyshev escape]\label{thm:chebyshev-escape}
Let $h,X\in\N$.  If
\[
 \theta(h)+m(A)<\theta(X),
\]
then there exists a prime $p$ satisfying
\[
 h<p\le X,
 \qquad p\notin A.
\]
\end{theorem}

\begin{proof}
Suppose no such prime exists.  Then every prime at most $X$ belongs either to
the set of primes at most $h$ or to $A$.  Since every logarithmic summand is
nonnegative,
\[
 \theta(X)\le\theta(h)+m(A),
\]
contradicting the strict hypothesis.
\end{proof}

Using explicit estimates of the form
\[
 \theta(h)\le (\log4)h
\]
and
\[
 X\log2-\log(X+1)-2\sqrt X\log X\le\theta(X),
\]
one obtains the following purely real sufficient condition:
\[
 (\log4)h+m(A)
 <X\log2-\log(X+1)-2\sqrt X\log X.
\]
The remaining analytic step is to choose an explicit $X=X(h,m(A))$ satisfying
this inequality.  This route is retained independently of the asymptotic
route below.

\section{Unbounded-prime route via relative source defects}

Let $P$ be an $abc$ point and let $B_n$ be genuine public IUT IV sources.
Write
\[
 \Corr_n(P)=\frac{20d_{\mathrm{mod},n}(P)}{\ell_n}.
\]
Suppose there are scalar sequences
$\gamma_n,\alpha_n,\beta_n,D_n,E_n$ such that
\begin{align*}
 \Corr_n(P)&\le\gamma_n,\\
 \Diff_n(P)&\le\alpha_n\Cond(P)+D_n,\\
 \Err_n(P)&\le\beta_n\Cond(P)+E_n.
\end{align*}
The constants $D_n,E_n$ may depend on $n$, because the final $abc$ constant
may depend on $\eps$.

\begin{theorem}[Relative absorption criterion]\label{thm:relative-absorption}
Assume that for every $\rho>0$ there exists a common $n$ satisfying
\[
 \gamma_n\le\rho,
 \qquad
 \alpha_n\le\rho,
 \qquad
 20\beta_n\le\rho.
\]
Then the logarithmic $abc$ conjecture holds.
\end{theorem}

\begin{proof}
The public IUT IV bound has the form
\[
 H(P)\le
 (1+\Corr_n(P))
 \bigl(\Diff_n(P)+\Cond_{\mathrm{Frey}}(P)\bigr)
 +20\Err_n(P)+C_0.
\]
The verified Frey conductor comparison gives
\[
 \Cond_{\mathrm{Frey}}(P)
 \le\Cond(P)+\log16.
\]
Substituting the three source bounds and collecting conductor terms yields
\[
 H(P)\le
 \bigl((1+\gamma_n)(1+\alpha_n)+20\beta_n\bigr)\Cond(P)
 +C_n.
\]
For $0\le\rho\le1$ and the common index above,
\[
 (1+\gamma_n)(1+\alpha_n)+20\beta_n
 \le (1+\rho)^2+\rho
 \le1+4\rho.
\]
Given $\eps>0$, choose $\rho=\min\{1,\eps/4\}$.  Then
$1+4\rho\le1+\eps$, and the additive terms form the allowed constant
$C_{\eps}$.
\end{proof}

\begin{theorem}[Uniform convergence criterion]\label{thm:uniform-convergence}
The common-index hypothesis of \cref{thm:relative-absorption} follows if
\[
 \gamma_n\to0,
 \qquad
 \alpha_n\to0,
 \qquad
 20\beta_n\to0.
\]
Consequently these three uniform limits imply $abc$.
\end{theorem}

\begin{proof}
For a fixed $\rho>0$, convergence of each sequence gives an eventual set of
indices on which the corresponding value is less than $\rho$.  The
intersection of three eventual sets is eventual and therefore nonempty.
Choose one common index in this intersection and apply
\cref{thm:relative-absorption}.
\end{proof}

\section{Necessity of simultaneous control}

The simultaneous nature of \cref{thm:relative-absorption} is not cosmetic.
Each defect can independently destroy a target coefficient $1+\eps$.

\begin{counterexample}[Persistent correction]
If $\Corr_n(P)\ge\delta>0$ uniformly while all other defects vanish, the
coefficient is at least $1+\delta$.  Hence no conclusion with
$\eps<\delta$ follows.
\end{counterexample}

\begin{counterexample}[Persistent different slope]
If $\alpha_n\ge\delta>0$ uniformly and the correction and error slope vanish,
the coefficient is at least $1+\delta$.  Again, no conclusion with
$\eps<\delta$ follows.
\end{counterexample}

\begin{counterexample}[Persistent error slope]
If $20\beta_n\ge\delta>0$ uniformly and the remaining defects vanish, the
coefficient is at least $1+\delta$.
\end{counterexample}

These examples do not rule out cancellation in a more refined source formula.
They rule out the proposed implication from bounds that leave one positive
linear defect uncontrolled.

\section{Two retained quantitative strategies}

The bounded and unbounded routes solve different problems and neither is
eliminated.

\begin{center}
\begin{tabularx}{\textwidth}{@{}lXX@{}}
\toprule
Route & Required theorem & Advantage \\
\midrule
Bounded Chebyshev route & An explicit $X$ satisfying the Chebyshev mass
inequality, together with a quantitative exceptional-prime bound & Direct
control of the first usable prime and every prime-dependent additive term \\
\addlinespace
Relative asymptotic route & Uniform limits
$\gamma_n,\alpha_n,20\beta_n\to0$ along genuine sources & Uses only
arbitrarily large admissible primes; no independent upper bound is needed \\
\bottomrule
\end{tabularx}
\end{center}

A future proof may combine them: use a bounded interval to control a finite
set of difficult local factors, while using relative absorption for all terms
that are sublinear in the conductor.

\section{Exact remaining geometric and analytic theorems}

The unresolved obligations can now be stated without broad interface labels.

\subsection{Fiber theorem}
One must construct, for the actual Frey and admissible-prime data:
\begin{enumerate}[label=(F\arabic*)]
\item the Tate-quotient theta-root cover obtained by globalizing
\cref{thm:kummer-fiber};
\item descent under $u\mapsto qu$ and treatment of the theta zero orbit;
\item the required orbicurve type and core;
\item the global and localized cartesian squares;
\item a valuation section formulated up to the correct equivalence of
normalized places;
\item the tempered fundamental group and its comparison with the profinite
\'etale group;
\item the canonical graph cusp and its compatibility with the global cusp.
\end{enumerate}

\subsection{Theorem 3.11 source theorem}
One must prove:
\begin{enumerate}[label=(S\arabic*)]
\item genuine Hodge-theater ordinary branches map to the native packet;
\item Ind1 and Ind2 are represented by the proposed packet maps;
\item the actual Ind3 target satisfies the relational log-link containment;
\item the three maps satisfy the pointwise norm estimates consumed by
\cref{thm:norm-reduction};
\item the possible-image union admits the intended finite-positive measured
hull or outer-measure interpretation;
\item the native procession volume is exactly $-|\log q|$;
\item the explicit envelope volume agrees with the public component formula
used in IUT IV.
\end{enumerate}

\subsection{Quantifier-correct final source family}
One must prove either:
\begin{enumerate}[label=(Q\arabic*)]
\item a bounded-prime theorem strong enough to select a usable prime below the
source-dependent upper threshold; or
\item the three uniform limits in \cref{thm:uniform-convergence}; or
\item another theorem whose algebra implies the same $1+\eps$ coefficient.
\end{enumerate}

\section{Lean correspondence}

The principal formal files are organized as follows.

\begin{center}
\begin{tabularx}{\textwidth}{@{}lX@{}}
\toprule
Mathematical result & Lean module \\
\midrule
Odd support and Frey signature &
\texttt{FreyOddPrimeSupport}, \texttt{FreyOddMultiplicativeCriterion} \\
Fixed global-core assembly &
\texttt{FreyTwelveTorsionGlobalCore} \\
Automatic quotient/cusp/decomposition choices &
\texttt{FiberAutomaticChoices} \\
Local multiplicative constructor &
\texttt{LocalMultiplicativeReductionBridge} \\
Unit-$c_4$ minimality &
\texttt{MinimalModelFromUnitC4} \\
Pointwise theta Kummer fiber &
\texttt{TateThetaKummerFiber} \\
Relational Ind3 source &
\texttt{UpperSemicompatiblePossibleImageSystem} \\
Pointwise norm reduction &
\texttt{Ind3NormEnvelopeReduction},
\texttt{NormControlledSourceGenerators} \\
Chebyshev bounded-prime criterion &
\texttt{ChebyshevPrimeEscape} \\
Relative defect absorption &
\texttt{RelativeSourceTermAbsorption} \\
Uniform convergence criterion &
\texttt{AsymptoticRelativeSourceTermAbsorption} \\
Necessity counterexamples &
\texttt{RelativeSourceTermNecessity},
\texttt{QuantitativeLargeImagePrimeNecessity} \\
\bottomrule
\end{tabularx}
\end{center}

A file is merged into \texttt{main} only after the repository's pinned Lean
kernel build succeeds.  Research branches preserve independent approaches,
and obsolete branches are closed only after their mathematical content has
been retained in a verified replacement or after a strict counterexample has
invalidated the proposed implication.

\section{Conclusion}

The three decisive obstacles have been reduced to substantially sharper
statements.

\begin{enumerate}[label=(\arabic*)]
\item The actual fiber is no longer an undifferentiated structure: its
automatic fields have been removed, and its first genuine local algebraic
piece is the separable Tate theta Kummer fiber.
\item The genuine Theorem~3.11 source need not satisfy an artificial exact
normal form.  A relational Ind3 presentation is sufficient, and its global
envelope theorem follows from pointwise norm control.
\item The auxiliary-prime problem has two rigorous routes.  The bounded route
is reduced to a Chebyshev inequality; the unbounded route is reduced to three
uniform relative-defect limits.  Counterexamples show why simultaneous
control is necessary.
\end{enumerate}

The remaining gap is difficult but explicit: globalize the theta-root cover
and construct the anabelian/tempered geometry; prove the genuine Hodge-theater
source maps and norm/volume calibration; and establish either quantitative
prime control or uniform asymptotic source decay.  Once these are proved, the
already formalized non-circular downstream chain yields
\cref{thm:abc-target}.  Until then, the correct status is a verified research
programme with exact open theorems, not a completed proof of $abc$.

\end{document}
